% !TeX program = xelatex
% BEGIN ORIGINAL: source/ru/main.tex
\documentclass[10pt,a4paper]{article}
\usepackage{fontspec}
\setmainfont{FreeSerif}
\setmonofont[Scale=0.85]{DejaVu Sans Mono}
\usepackage[english,russian]{babel}
\usepackage[margin=22mm]{geometry}
\usepackage{amsmath,amssymb,amsthm,mathtools}
\usepackage{booktabs,tabularx}
\usepackage{tikz}
\usetikzlibrary{positioning}
\usepackage{microtype}
\usepackage[hidelinks]{hyperref}
\hypersetup{pdftitle={Контрактно-относительная теория преобразований: от сертифицированного открытия к замыканию метаполитик},pdfauthor={Александр Странд},pdfsubject={Бюджетное открытие и сертифицированная замена конечных метаполитик}}
\newtheorem{theorem}{Теорема}[section]
\newcommand{\Nat}{\mathbb N}
\newcommand{\CanonRef}{\mathcal D}
\newcommand{\bill}{\operatorname{bill}}
\newcommand{\spent}{\operatorname{spent}}
\newcommand{\funds}{\operatorname{funds}}
\newcommand{\reserve}{\operatorname{reserve}}
\newcommand{\balance}{\operatorname{balance}}
\newcommand{\PB}{\operatorname{PB}}
\newcommand{\Pres}{\operatorname{Pres}}
\newcommand{\Cov}{\operatorname{Cov}}
\newcommand{\finish}{\operatorname{finish}}
\newcommand{\Correct}{\operatorname{Correct}}
\newcommand{\rc}[1]{\allowbreak\texttt{#1}\allowbreak}
\setlength{\emergencystretch}{3em}
\tolerance=2500
\begin{document}
% BEGIN ORIGINAL: source/ru/sections12.tex
% Русская версия RC4. Исходная структура, метки и формулы сохранены.
\title{\textbf{Контрактно-относительная теория преобразований:}\\
\large От сертифицированного открытия к замыканию метаполитик}
\author{Александр Странд (Aleksandr Strand)\\[-1pt]{\small\href{mailto:strandlex@gmail.com}{strandlex@gmail.com}}}
\date{\vspace{-11pt}Ноябрь 2026 г.}
\maketitle

\begin{abstract}
Представлена \textbf{контрактно-относительная теория преобразований (CRTT)} --- механизированная система бюджетного открытия и сертифицированного изменения исследовательских политик. Исходя из достаточности информации относительно контракта, мы разрабатываем конструктивный синтез исследований, сертифицированный перенос методов, возобновляемый поиск морфизмов, компиляцию препятствий и платное приобретение свидетельств, достаточных для решения. Эти конструкции обеспечивают механизм замены с сохранением продолжения, отделяющий семантическую корректность от экономического выигрыша и ликвидности.

Главный результат устанавливает композиционное замыкание конечных корректно типизированных расписаний обновления метаполитик при фиксированных исходной и целевой задачах, фиксированном поставщике свидетельств и финансируемом начальном продолжении. Исполнимая реализация сохраняет членство в каноническом результате, точную характеризацию невозможности, баланс ресурсов и предшествующую сертифицированную историю. Принятые обновления строго улучшают локальное значение полной модельной стоимости; кумулятивная монотонность гарантируется для трасс без отклонённых обновлений. Контрпримеры показывают, почему корректное обучение и локально выгодные замены не обязательно дают глобальный экономический выигрыш.

Каноническая реализация включает 21 исходный файл Lean и 542 аудируемые декларации в Lean~4.36.0, независимые конечные перекрёстные проверки и девять исполнимых свидетелей релиза. Гарантии ограничены объявленной конечной детерминированной моделью; произвольное самопереписывание, автономное порождение свидетельств и улучшение физической производительности не предполагаются.
\end{abstract}

\section{Введение и происхождение исследования}
\label{sec:intro}
\subsection{От сертифицированной функциональной замены к препятствиям CRTT}
Исследовательская программа CRTT возникла из задачи \emph{сертифицированной функциональной замены} (Certified Functional Replacement, CFR): когда один внутренний механизм сложной вычислительной или нейронной системы можно заменить другим, сохранив предписанное контрактом поведение? Эксперименты CFR мотивировали различение трёх возможных барьеров замены:
\begin{enumerate}
    \item \textbf{Операторный барьер:} в классе кандидатов нет допустимого оператора, реализующего требуемое поведение.
    \item \textbf{Репрезентационный барьер:} у слоя наблюдения может отсутствовать общее допустимое действие, даже если каждое его состояние разрешимо по отдельности.
    \item \textbf{Барьер выбора:} даже при достаточном наблюдении механизм выбора может не находить допустимый оператор. Эмпирические сбои калибровки мотивируют такое различение, но сами по себе не доказывают репрезентационную невозможность.
\end{enumerate}

Формализация этих барьеров в бюджетной системе с проверяемыми доказательствами выявила три новых класса проблем. Они не являются попарными эквивалентами барьеров CFR, а возникают при переходе от статической замены к работающим исследовательским системам: \emph{структурная непереносимость} между задачами, \emph{затраты на приобретение и поиск} методов и свидетельств, \emph{ограничения миграции и ликвидности} при замене действующего контроллера. Отсюда главный исследовательский вопрос CRTT:
\begin{quote}
\emph{При каких условиях исследовательская система может изменять собственные методы и управляющие политики, сохраняя семантическую корректность, конечные ресурсы и происхождение результатов, подтверждённое доказательствами?}
\end{quote}

\subsection{Архитектура как последовательность переходов}
CRTT отвечает на этот вопрос посредством шести качественных математических переходов (рис.~\ref{fig:arch-map}). Каждый слой разрешает конкретное ограничение предыдущего (табл.~\ref{tab:genesis}). В каноническом стеке стоимости, остатки и резервы --- натуральные учётные величины объявленной детерминированной модели, а не физическое время или память.

\begin{figure}[h]
\centering
\small
\begin{tikzpicture}[
    node distance=0.35cm and 0.45cm,
    box/.style={rectangle, draw=black!75, rounded corners=2pt, align=center,
                text width=3.9cm, minimum height=1.25cm, fill=black!3},
    arrow/.style={->, >=stealth, thick}
]
\node[box] (b1) {\textbf{Раздел 2}\\Контрактная достаточность\\\emph{Что нужно знать?}};
\node[box, right=of b1] (b2) {\textbf{Раздел 3}\\Конструктивное исследование\\\emph{Как получить знание?}};
\node[box, right=of b2] (b3) {\textbf{Раздел 4}\\Компиляция препятствий\\\emph{Как превратить сбой в правило?}};
\node[box, below=0.55cm of b1] (b4) {\textbf{Раздел 5}\\Метасвидетельства для решения\\\emph{Как выбирать методы?}};
\node[box, right=of b4] (b5) {\textbf{Раздел 6}\\Сертифицированная замена\\\emph{Как менять контроллер?}};
\node[box, right=of b5] (b6) {\textbf{Раздел 7}\\Композиционное замыкание\\\emph{Как сохранить гарантии?}};
\draw[arrow] (b1) -- (b2);
\draw[arrow] (b2) -- (b3);
\draw[arrow] (b3.south) -- ++(0,-0.25) -| (b4.north);
\draw[arrow] (b4) -- (b5);
\draw[arrow] (b5) -- (b6);
\end{tikzpicture}
\caption{Концептуальное развитие CRTT: стрелки означают переход между исследовательскими вопросами, а не обязательно импорт исходников.}
\label{fig:arch-map}
\end{figure}

\begin{table}[h]
\centering
\footnotesize
\begin{tabular}{@{}>{\raggedright\arraybackslash}p{3.6cm}>{\raggedright\arraybackslash}p{5.2cm}>{\raggedright\arraybackslash}p{4.6cm}@{}}
\toprule
\textbf{Установленный результат} & \textbf{Обнаруженная граница} & \textbf{Следующий объект исследования} \\
\midrule
Функциональная заменимость & Допустимого оператора может не быть & Препятствие \\
Достаточность наблюдения & Требуемая информация заранее неизвестна & Стратегия наблюдения \\
Оптимальная стратегия & Повторный синтез в новой задаче дорог & Перенос \\
Переносимость & Требуется найти морфизм переноса & Конструктивное открытие \\
Полное открытие & Исчерпывающий поиск слишком дорог & Сертифицированное отсечение \\
Обучение правилам & Корректное обучение может не окупаться & Метаконтроллер \\
Выбор контроллера & Точные стоимости методов известны частично & Метасвидетельства \\
Достаточные метасвидетельства & Замена может нарушить продолжение & Сертифицированная миграция \\
Корректная миграция & Повторные замены взаимодействуют в истории & Композиционные законы трасс \\
Замена контроллера & Сама политика выбора может быть недостаточной & Замыкание метаполитик \\
\bottomrule
\end{tabular}
\caption{Исследовательские вопросы, обусловленные границами предшествующих конструкций.}
\label{tab:genesis}
\end{table}

\subsection{Путеводитель и дополнительный архив}
Разделы~\ref{sec:foundations}--\ref{sec:meta-evidence} строят конструктивный и экономический аппарат для формулировки и доказательства главных результатов разделов~6 и~7. Полные исходники Lean~4.36.0 (21 канонический файл, 542 аудируемые декларации), исторические записи этапов и девять свидетелей исполнения CLI содержатся в согласованном архиве релиза (\texttt{CRTT-v1.0-FULL-ARCHITECTURE-RECONCILED.zip}), служащем формальным приложением.

\section{Контрактно-относительные основания}
\label{sec:foundations}
\subsection{Контракты, достаточность и критерий существования}
Зафиксируем пространство состояний $X$, непустое пространство действий $A \neq \varnothing$ и предикат контракта $V : X \times A \to \mathrm{Prop}$. Для каждого $x \in X$ определим множество допустимых действий $A_C(x) = \{ a \in A \mid V(x, a) \}$. Отображение наблюдения $\phi : X \to Z$ задаёт слой $\phi^{-1}(z) = \{ x \in X \mid \phi(x) = z \}$ над каждой наблюдаемой сигнатурой $z \in \phi(X)$. Тотальная политика $\pi : Z \to A$ \emph{допустима относительно контракта через $\phi$}, если $V(x, \pi(\phi(x)))$ выполняется для всех $x \in X$.

\begin{theorem}[Критерий контрактно-относительной достаточности (\texttt{CRTT.selector\_iff})]
\label{thm:sufficiency}
Пусть $A \neq \varnothing$. Отображение наблюдения $\phi : X \to Z$ допускает тотальную контрактно-допустимую политику $\pi : Z \to A$ тогда и только тогда, когда каждый реализуемый слой имеет общее допустимое действие:
\begin{equation}
\label{eq:sufficiency}
\exists \pi : Z \to A,\; \forall x \in X,\; V(x, \pi(\phi(x)))
\iff
\forall z \in \phi(X),\; \bigcap_{\phi(x) = z} A_C(x) \neq \varnothing.
\end{equation}
\end{theorem}
\begin{proof}
($\Rightarrow$) Если $\pi$ удовлетворяет $V(x, \pi(\phi(x)))$ для всех $x$, то для любого $z \in \phi(X)$, $\pi(z) \in \bigcap_{\phi(x)=z} A_C(x)$.
($\Leftarrow$) Для каждого $z \in \phi(X)$ выберем свидетель $a_z \in \bigcap_{\phi(x)=z} A_C(x)$ (в общем случае с использованием классического выбора). Для $z \notin \phi(X)$ положим $\pi(z) = a_0$, где $a_0 \in A$ существует в силу $A \neq \varnothing$. Полученная тотальная политика удовлетворяет $V(x, \pi(\phi(x)))$ для всех $x \in X$.
\end{proof}
Без $A \neq \varnothing$ критерий неверен при $X = \varnothing$ и $Z \neq \varnothing$. Кроме того, этот ранний общий результат использует классический выбор, в отличие от конструктивного конечного поиска канонического стека, аудируемые декларации которого не зависят от \texttt{Classical.choice}.

\subsection{Коллективные препятствия и грубейшие разбиения}
Из критерия~\eqref{eq:sufficiency} следует, что достаточность нарушается, если некоторый слой имеет пустое пересечение допустимых действий. Важно, что уточнение $\phi$ не всегда устраняет препятствие: если $A_C(x) = \varnothing$ для некоторого $x \in X$, даже полное наблюдение состояния не даёт допустимого действия.

Даже если каждое состояние разрешимо ($A_C(x) \neq \varnothing$), препятствия не обязательно попарны. Определим совместимость миров:
\[
x \sim_C y \iff A_C(x) \cap A_C(y) \neq \varnothing.
\]
Рассмотрим три мира с множествами допустимых действий $A_C(x_1)=\{a,b\}$, $A_C(x_2)=\{b,c\}$ и $A_C(x_3)=\{a,c\}$ (\texttt{SufficientPartitions.triangle\_compat\_all}, \texttt{triangle\_no\_common\_action}). Любые два мира имеют общее действие, но $\bigcap_{i=1}^3 A_C(x_i) = \varnothing$: препятствие строго коллективное. Это явление ограничивает существование наибольшего достаточного разбиения в порядке уточнения (а не кодировки минимального размера).

\begin{theorem}[Существование грубейшего достаточного разбиения]
\label{thm:coarsest}
Грубейшее достаточное разбиение $X$ существует тогда и только тогда, когда $\sim_C$ является отношением эквивалентности и разбиение, порождённое $\sim_C$, само достаточно.
\end{theorem}
\begin{proof}[Эскиз доказательства]
Если $\sim_C$ --- достаточное отношение эквивалентности, любое достаточное разбиение обязано разделять несовместимые миры и, значит, уточняет $\sim_C$; следовательно, $\sim_C$ является грубейшим. Обратно, пусть грубейшее достаточное разбиение $\mathcal{P}^*$ существует. Достаточность $\mathcal{P}^*$ влечёт $A_C(w)\ne\varnothing$ для каждого $w$. Тогда любую совместимую пару $x\sim_C y$ можно объединить, оставив остальные состояния одноэлементными блоками. Поскольку любое достаточное разбиение уточняет $\mathcal{P}^*$, $x$ и $y$ лежат в одном классе. С другой стороны, достаточность $\mathcal{P}^*$ требует совместимости любых двух состояний одного класса. Следовательно, $\sim_C$ совпадает с отношением эквивалентности $\mathcal{P}^*$ и достаточно.
\end{proof}
Соответствующая декларация --- \texttt{coarsest\_exists\_iff} в \texttt{SufficientPartitions.lean}. В примере с треугольником $\sim_C$ --- универсальное отношение (а значит, эквивалентность), но оно недостаточно; грубейшего достаточного разбиения нет. Напротив, для точной цели $g:X\to Y$ равенство целевых значений само является грубейшим достаточным отношением (\texttt{exact\_target\_coarsest}).

\newpage
\subsection{Экономические основания: накопление библиотеки и минимаксное сожаление}
Установив достаточность, исследуем, когда накопление методов или приобретение наблюдений окупается при конечных бюджетах:
\begin{itemize}
    \item \textbf{Идеализированное накопление и строгий выигрыш.} Для фиксированного предиката допустимой трассы, если повторное использование не дороже открытия и дополнительная плата за поиск отсутствует, бюджетная достижимость монотонна по включению библиотек. Обозначим через $C_{\mathrm{ind}}^*$ и $C_{\mathrm{cum}}^*$ достигаемые минимумы по одним и тем же допустимым планам, причём $C_{\mathrm{ind}}(p)=C_{\mathrm{cum}}(p)+S(p)$ для каждого допустимого $p$. Строгий кумулятивный выигрыш имеет место тогда и только тогда, когда некоторый план $p$ даёт сертифицированную экономию $S(p)$, превышающую разрыв до независимого оптимума:
    \begin{equation}
    \label{eq:t15}
    C_{\mathrm{cum}}^* < C_{\mathrm{ind}}^* \iff \exists p :\; S(p) > C_{\mathrm{ind}}(p) - C_{\mathrm{ind}}^*.
    \end{equation}
    \item \textbf{Обязательная плата за поиск и онлайн-пороги.} Обязательная плата за поиск может нарушать монотонность: расширение библиотеки увеличивает стоимость с $2$ до $3$ при бюджете $2$. В модели скрытого метода неблагоприятно выбранный повторяющийся скрытый метод даёт сожаление $D-R$ для детерминированного первого выбора, когда цена открытия равна $D$, повторного использования --- $R\le D$ и имеется не менее двух задач. В отдельной пороговой модели повторяющейся задачи с ценой аренды $L$ и $R<L\le D$ стратегия точки окупаемости $2$-конкурентна относительно офлайн-оптимума.
    \item \textbf{Минимаксное сожаление и ценность наблюдения.} Для конечного списка сценариев, непустого конечного списка действий и разрешимого равенства наблюдений пусть $c(a,x)\in\mathbb N$. Определим сожаление действия $r(a,x) = c(a,x) - \min_{b \in A} c(b,x)$ и минимаксное сожаление наблюдения $\phi$ в худшем случае:
    \begin{equation}
    \label{eq:t18-regret}
    R^*(\phi) = \max_{z \in \phi(X)} \min_{a \in A} \max_{x : \phi(x)=z} r(a,x).
    \end{equation}
    Максимум пустого списка сценариев равен нулю. Минимаксный критерий характеризует $\varepsilon$-достаточность существованием в каждом реализуемом слое общего действия с сожалением не более $\varepsilon$. В примере XOR отсутствие наблюдения или наблюдение любого одного бита даёт сожаление $1$, обоих битов --- $0$. Приобретение фиксированного уточнения $\phi'$ за цену $c_{\mathrm{obs}}$ строго выгодно в точности тогда, когда
    \begin{equation}
    \label{eq:t18-acq}
    c_{\mathrm{obs}} + R^*(\phi') < R^*(\phi).
    \end{equation}
\end{itemize}

\subsection{От достаточности к исследованию}
Эти результаты включают онлайн-пороговую модель, но ещё не дают общего синтезатора деревьев ограниченной глубины и типизированного механизма переноса задач, разрабатываемых ниже. Это мотивирует слой конструктивного исследования и переноса (раздел~3).

Результаты оснований имеют отдельные исторические доказательные артефакты. Самодостаточный канонический стек устанавливает теорему релиза независимо от этих ветвей; их происхождение описано в разделе~8.
% END ORIGINAL: source/ru/sections12.tex
\clearpage
% BEGIN ORIGINAL: source/ru/section3.tex
\section{Конструктивное исследование и перенос}\label{sec:investigation}
Достаточность задаёт критерий для фиксированного наблюдения, но не стратегию приобретения информации. Сначала синтезируем конечные адаптивные исследования, затем выясним, когда синтезированную стратегию можно перенести на другую задачу. Объявленные списки миров, действий и портов конечны; равенство портов разрешимо. Оценка эффективности исчерпывающего синтеза не утверждается.

\subsection{Синтез адаптивного исследования}
Задача задаётся списком миров $E$, непустым списком действий $\mathcal A$, списком портов $Q$, функцией чтения $r(q,w)\in\{0,1\}$, ценами портов $c(q)\in\Nat$, терминальными потерями $L(a,w)\in\Nat$ и пределом глубины $d$. Исследование --- конечное дерево
\[
 t ::= \operatorname{act}(a)\mid\operatorname{observe}(q,t_0,t_1).
\]
Допустимость требует объявленных действий и портов, глубины не более $d$ и отсутствия повторного порта вдоль пути (все вхождения приобретённого порта удаляются из оставшегося списка). Для реализованной стоимости $J_w(t)$ семантика такова: $J_w(\operatorname{act}(a))=L(a,w)$ и $J_w(\operatorname{observe}(q,t_0,t_1))=c(q)+J_w(t_{r(q,w)})$. Оплачивается только реализованная ветвь. Пусть $W_E(t)=\max_{w\in E}J_w(t)$, максимум пустого списка равен нулю, а $V_d(E,Q)$ --- оптимальная худшая стоимость допустимого дерева. Обозначим через $E_b(q)$ фильтрацию $E$ по значению чтения $b$, и $L_*(E)=\min_{a\in\mathcal A}\max_{w\in E}L(a,w)$.

\begin{theorem}[Оптимальное конечное исследование]\label{thm:investigation}
Исполнимый решатель достигает $V_d(E,Q)$ допустимым деревом. При $E\ne\varnothing$,
\begin{align}
 V_0(E,Q)&=L_*(E),\\
 V_{d+1}(E,Q)&=\min\Bigl(\{L_*(E)\}\cup\{c(q)+\max_{b\in\{0,1\}} \\
 &\hspace{24mm} V_d(E_b(q),Q\setminus\{q\}):q\in Q\}\Bigr).
\end{align}
Для пустого $E$ значение равно нулю. Для жёсткого допуска терминальных потерь $\epsilon$ и бюджета наблюдений $b$ заменим терминальную потерю на $0$, если $L(a,w)\le\epsilon$, и на $b+1$ иначе. Оптимум такой ограниченной модели не превосходит $b$ тогда и только тогда, когда допустимое дерево соблюдает обе границы во всех объявленных мирах.
\end{theorem}
\begin{proof}
Индукция по глубине. Любое допустимое дерево либо является терминальным действием, либо начинается с объявленного порта и двух допустимых поддеревьев меньшей глубины. Для непустого $E$ худшая стоимость узла наблюдения равна плате плюс максимум стоимости ветвей: одна из ветвей реализуется, и все стоимости неотрицательны. Решатель выбирает наименее дорогого кандидата среди оптимального терминального действия и рекурсивно синтезированных узлов наблюдения. По индукции стоимость любого конкурирующего дерева не меньше минимума, а выбранный кандидат допустим. В ограниченной модели нарушающий допуск лист стоит не менее $b+1$, допустимый даёт нуль; поэтому ограничение худшей стоимости в точности проверяет жёсткую достижимость. См. \rc{solve\_legal}, \rc{solve\_optimal}, \rc{bellman} и \rc{epsilon\_budget\_iff} в \rc{AdaptiveObservation.lean}.
\end{proof}
Адаптивный оптимум не превосходит лучшего допустимого фиксированного расписания (\rc{adaptive\_le\_static}). В свидетеле-мультиплексоре читается бит селектора, затем указанный бит данных: стоимость $2$ вместо статического оптимума $3$ при единичных ценах и потере $4$ за неверный ответ. Свидетель XOR также исключает общее жадное правило одношаговой ценности. Синтез точен в объявленном языке деревьев, а не по произвольным программам.

\subsection{Перенос и экономика библиотеки}
Для задач $S,T$ морфизм переноса $M=(h,p,a)$ отображает целевые миры в исходные, исходные порты в целевые, исходные действия в целевые, сохраняя объявленное членство. Карта портов инъективна на объявленных исходных портах, $d_S\le d_T$, и для всех объявленных аргументов
\begin{align*}
 r_T(p(q),w)&=r_S(q,h(w)), & c_T(p(q))&\le c_S(q),\\
 L_T(a(u),w)&\le L_S(u,h(w)).
\end{align*}
Перенос $M_*t$ --- синтаксическая подстановка меток портов и действий, не зависящая от неизвестного целевого мира.
\begin{theorem}[Сертифицированный перенос]\label{thm:transfer}
Для любого допустимого в исходной задаче дерева $t$ дерево $M_*t$ допустимо в целевой задаче, и $J^T_w(M_*t)\le J^S_{h(w)}(t)$ для каждого объявленного целевого мира. Плата за наблюдения и терминальные потери не возрастают по отдельности. Следовательно, $V_T\le V_S$; взаимные морфизмы влекут равенство оптимумов. Тождественность и композиция морфизмов индуцируют тождественность и композицию переносов деревьев.
\end{theorem}
\begin{proof}
Структурная индукция по дереву даёт совпадение выбора ветвей по равенству чтений и неравенства стоимости по сложению. Инъективность гарантирует, что удаление наблюдённого порта не удалит образ другого исходного порта. Неравенство глубин обеспечивает допустимость. Максимизация по целевым мирам и применение к исходному оптимизатору доказывают неравенство значений. Композиция есть композиция функций (контравариантная по мирам), а подстановки деревьев компонуются. Это результаты \rc{transfer\_legal}, \rc{transfer\_runCost\_le}, \rc{transfer\_hard\_contract}, \rc{optimum\_le}, \rc{mutual\_optimum\_eq} и \rc{transfer\_compose} в \rc{InvestigationTransfer.lean}.
\end{proof}
Равенство оптимумов само по себе не влечёт существования морфизма; конечные свидетели равных оптимумов демонстрируют нарушение инъективности и сохранения потерь. Перенос также требует затрат. Для указанного маршрута с накладными расходами $k$ и стоимостью целевого исполнения $e\ge V_T$ повторное использование строго дешевле нового синтеза за $D$ с оптимальным исполнением тогда и только тогда, когда $k+(e-V_T)<D$. Явный селектор библиотеки минимизирует стоимость по новому синтезу и перечисленным сертифицированным предложениям; при цене поиска $s$ и цене нового синтеза с исполнением $F$ поиск окупается тогда и только тогда, когда некоторое найденное предложение удовлетворяет $s+k+e<F$. Полнота относительно семантического оракула предложений требует полного по стоимости допущенного списка, а не просто корректного подмножества. Это \rc{paid\_candidate\_iff}, \rc{search\_profitable\_iff} и \rc{complete\_search\_matches\_oracle} в \rc{InvestigationLibrary.lean}.

\subsection{Конструктивное обнаружение морфизмов}
Теорема переноса полезна, лишь когда доступен её свидетель. Для индексированных конечных задач (объявлены все действия, включая действие по умолчанию) перечислим все карты $p:Q_S\to Q_T$ и $a:\mathcal A_S\to\mathcal A_T$. Для фиксированной пары определим
\[
 \mathcal Y_w(p,a)=\{y\in X_S: \forall q\in Q_S,\ r_T(p(q),w)=r_S(q,y),\quad
 \forall u\in\mathcal A_S,\ L_T(a(u),w)\le L_S(u,y)\}.
\]
\begin{theorem}[Каноническое обнаружение]\label{thm:discovery}
Фиксированная пара $(p,a)$ продолжается до морфизма тогда и только тогда, когда соблюдены условия глубины и портов и $\mathcal Y_w(p,a)\ne\varnothing$ для каждого целевого мира $w$. Выбор первого свидетеля в каждой конечной строке исходных миров строит морфизм. Список $\CanonRef=\rc{discoverAll}(S,T)$ содержит представителя с теми же картами портов и действий для каждого допустимого морфизма; он пуст тогда и только тогда, когда морфизмов нет.
\end{theorem}
\begin{proof}
Допустимая карта миров даёт элемент каждой строки. Обратно, каждому целевому миру можно независимо назначить первый элемент строки: морфизм не накладывает межмировых ограничений. Тем самым выполнены все обязательства чтения и потерь; отдельно проверенные условия глубины, цены и инъективности не меняются. Полный перебор карт включает пару любого допустимого морфизма. Непустой результат сам содержит элемент типа морфизмов. См. \rc{fixed\_pair\_iff}, \rc{discover\_complete} и \rc{empty\_iff\_no\_morphism} в \rc{MorphismDiscovery.lean}.
\end{proof}
Каноническому списку не требуется перечислять все возможные карты миров: перенесённое дерево зависит только от $p,a$. Это устраняет полный перебор глобальных карт миров, но не потенциально дорогой перебор остальных пар. Следующая задача --- сохранить эталонную семантику при прерывании, отсечении и обучении.
% END ORIGINAL: source/ru/section3.tex
\clearpage
% BEGIN ORIGINAL: source/ru/section4.tex
\section{Бюджетное открытие и компиляция препятствий}\label{sec:obstruction}
Конечный эталонный список $\CanonRef$ определяет, что должен восстановить поиск; его предварительное вычисление противоречило бы цели бюджетного поиска. Используем оплачиваемый исполнитель, корректность которого определяется относительно $\CanonRef$ без вызова эталонного перечисления на каждом шаге.

\subsection{Оплаченные префиксы и сертифицированное отсечение}
Детерминированная функция котировки сопоставляет состоянию $u$ отсутствие шага или пару $(\eta, k)$ с положительной натуральной платой $\eta$ и отложенным вычислением следующего состояния $k:\mathrm{Unit}\to U$. Счёт $(u,B,D)$ хранит остаток средств $B$ и зарегистрированные расходы $D$. Если $\eta\le B$, исполнитель вызывает $k$ и меняет счёт на $(k(),B-\eta,D+\eta)$; иначе останавливается, не вызывая $k$. Ограниченный запуск считает шаги отдельно от средств; запаса в $B$ шагов достаточно, чтобы достичь состояния без доступного по цене следующего шага. Индукция даёт точное сохранение $B'+D'=B+D$ и любого индуктивного инварианта котируемых переходов (\rc{accounting}, \rc{invariant}, \rc{budget\_bound} в \rc{BudgetedDiscovery.lean}).

\begin{theorem}[Точное бюджетное продолжение]\label{thm:resume}
Для фиксированной детерминированной функции котировки, не зависящей от бюджета и имеющей положительные платы, пусть $\mathsf E$ --- исполнитель, а $\mathsf U_b$ прибавляет $b$ к остатку средств. Тогда
\[
 \mathsf E\mathsf U_{b_1+b_2}(s)=\mathsf E\mathsf U_{b_2}(\mathsf E\mathsf U_{b_1}(s)).
\]
Это равенство включает состояние, остаток и расходы даже при циклической системе переходов.
\end{theorem}
\begin{proof}
Каждый префикс, доступный при меньшем бюджете, доступен и при большем; оплата коммутирует с добавлением разности. Применим индукцию по оплаченному префиксу, считая случаи недоступного и отсутствующего шага стационарными. Положительные платы ограничивают число оплаченных шагов остатком средств, откуда следует тождество для исполнителя до исчерпания доступных шагов. Это \rc{execute\_topUp\_add} в \rc{ResumeLaws.lean}. Теорема не влечёт семантического исчерпания поиска: цикл может просто израсходовать бюджет.
\end{proof}
В поиске морфизмов частичное назначение портов $\rho$ задаёт область продолжающих его морфизмов. Сертифицированное правило отсечения содержит булев тест $k(\rho)$, плату и свидетельство
\[
 \forall M,\rho,\quad \operatorname{Extends}(M,\rho)\Longrightarrow k(\rho)=\mathrm{true}.
\]
Ложный результат исключает всю область. Правила могут проверять инъективность, цены портов, глубину или совместимость чтений; корректные правила компонуются. Поиск с отсечением использует тот же конструктор канонической пары, что и раздел~3.

\begin{theorem}[Отсечение сохраняет каноническое членство]\label{thm:prune}
Для каждого сертифицированного правила и морфизма $x$ членство в полученном исчерпывающем списке с отсечением эквивалентно $x\in\CanonRef$.
\end{theorem}
\begin{proof}
Сохранённое полное назначение проверяется эталонным конструктором пары, что даёт корректность. Для полноты проследим префиксы карты портов любого канонического представителя; корректность каждого правила обеспечивает сохранение всех этих префиксов. Затем рассматривается эталонная карта действий и получается тот же канонический сертификат. Это \rc{prunedDiscover\_iff} в \rc{CertifiedPruning.lean}; равенство порядка и кратности списков не требуется.
\end{proof}
Корректное отсечение не обязано экономить ресурсы: при $n_1\le n_0$ посещённых узлах, единой цене узла $c$ и цене правила на узел $k$ оно окупается в точности при $n_1 k<(n_0-n_1)c$ (\rc{pruning\_profitable\_iff}). Корректное правило без эффекта повышает стоимость, если за него взимается плата. Более развитый исполнитель ниже отдельно учитывает фактические вызовы правил с досрочным прекращением проверки.

\subsection{Финансируемый фронтир поиска}
Фронтир $f=(J,A,\Lambda)$ хранит оставшиеся задания $J$, выданные канонические сертификаты $A$ и оплаченные поисковые квитанции $\Lambda$. Задание либо расширяет область назначений портов с оставшейся очередью правил, либо обрабатывает оставшийся суффикс карт действий. Корректный фронтир содержит только канонические сертификаты и покрывает каждый ещё не выданный канонический сертификат хотя бы одним оставшимся заданием; точный предикат дан в разделе~6.1. Разбиение задания, исключение сертифицированной области или проверка пары сохраняют этот предикат и дописывают новый сертификат, если он найден. Доказательство использует локальное покрытие и корректность, а не исполнение эталонного перечисления.

Пусть $P(C,f)$ --- натуральный потенциал оставшихся обязательств при фиксированной библиотеке правил и базовых тарифах. Бюджеты его заданий ограничивают полную работу ветвления, включая платы за правила. Каждый котируемый платный шаг удовлетворяет
\begin{equation}\label{eq:funded-potential}
 P(C,f')+\eta\le P(C,f),\qquad\eta>0.
\end{equation}
\begin{theorem}[Финансируемое исчерпание]\label{thm:funded-search}
Начиная с корректного фронтира при $B\ge P(C,f)$, исполнитель исчерпывает задания, сохраняет прежние сертификаты как префикс списка и возвращает список с членством в точности как в $\CanonRef$. Пустой список сертификатов при исчерпанном фронтире доказывает отсутствие морфизма.
\end{theorem}
\begin{proof}
Уравнение~\eqref{eq:funded-potential} сохраняет финансирование после платежа и строго уменьшает потенциал. При остановке наличие оставшегося задания означало бы доступный по цене положительный шаг, что противоречит остановке. Локальная корректность индуктивна, а пустой список заданий устраняет дизъюнкт оставшегося покрытия. Сохранение префикса следует по индукции из дописывания сертификатов. Это \rc{quote\_potential}, \rc{execute\_funded\_exhausts}, \rc{execute\_full}, \rc{execute\_certificate\_prefix} и \rc{exhausted\_impossible} в \rc{BudgetedPruning.lean}.
\end{proof}
Без предпосылки финансирования классификатор различает найденный свидетель, пустой исчерпанный фронтир с доказательством невозможности и неопределённый исход при пустом списке сертификатов и неисчерпанном фронтире. Возвращённый счёт допускает продолжение. Найденный до исчерпания свидетель не означает полноту перечисления.

\subsection{От пустой строки к переиспользуемому правилу}
Для целевого мира $w$ частичное назначение портов $\rho$ определяет строку совместимых чтений
\[
 Y_w(\rho)=\{y\in X_S: \forall(q,q')\in\rho,\ r_S(q,y)=r_T(q',w)\}.
\]
Если $Y_w(\rho)=\varnothing$, добавление ограничений не может восстановить элемент. Компилятор препятствий удаляет ограничения, сохраняя пустоту строки, и получает подмножество $K\subseteq\rho$, которое остаётся несовместимым, но перестаёт быть таковым при удалении любого из значений ограничений. Таким образом, $K$ неприводимо по включению, но не обязательно имеет минимальную мощность.
\begin{theorem}[Сертифицированная компиляция препятствий]\label{thm:core}
Для пустой строки чтений \rc{extract} строит такое неприводимое $K$ со свидетельством. Ни одно назначение, содержащее $K$, не допускает продолжающего морфизма. Скомпилированный тест отклоняет в точности назначения, содержащие $K$, и удовлетворяет контракту сертифицированного правила.
\end{theorem}
\begin{proof}
Процедура сокращения просматривает исходные ограничения, удаляя значение, если несовместимость сохраняется. Если удаление не удалось, монотонность несовместимости гарантирует, что последующие удаления не могут снова сделать результат этого удаления несовместимым. Поэтому оставшееся подмножество неприводимо. Если бы существовал морфизм, продолжающий назначение, содержащее $K$, его карта миров в $w$ дала бы элемент пустой строки. Булев тест включения точен по отражению членства. См. \rc{shrink\_spec}, \rc{extract}, \rc{core\_excludes}, \rc{compileCore\_rejects} в \rc{ImpasseCompilation.lean}.
\end{proof}
Перенос ядра на новую пару задач допускается при наличии карты ограничений, отображения новых исходных миров в старые и доказательства отражения чтений для каждого перенесённого ограничения (\rc{transportCore}); неприводимость и экономическое превосходство перенесённого правила не гарантируются.

Обучающийся исполнитель делает компиляцию эндогенной. Состояние поиска включает активную библиотеку ядер, исходные и сокращённые события обучения, ожидающую компиляцию и оплаченные квитанции. Поиск в библиотеке, зондирование, компиляция и базовое продвижение поиска --- отдельные переходы с положительными платами. Ядро вставляется с поглощением: если старое ядро уже исключает всю его область, новое не устанавливается; иначе удаляются ядра, ограничения которых содержат новое ядро. Итоговая исключаемая область в точности равна объединению прежней области и области нового ядра (\rc{insertCore\_rejects}); предыдущие события обучения сохраняются в истории. Теоремы \rc{execute\_correct}, \rc{execute\_trace}, \rc{execute\_split} и \rc{execute\_exhausted\_iff} в \rc{LearningDiscovery.lean} сохраняют семантические и учётные гарантии. Исчерпание условно и не гарантируется произвольным лимитом обучения.

Эта конструкция делает отрицательные результаты операционно переиспользуемыми. Она не устанавливает универсально более дешёвого обучающегося исполнителя. В объявленном конечном стенде обучение было дороже фиксированного контроллера правила чтения во всех 515 полностью завершённых случаях. Поэтому следующая задача --- выбор самого исследовательского механизма, а не безусловное обучение.
% END ORIGINAL: source/ru/section4.tex
\clearpage
% BEGIN ORIGINAL: source/ru/section5.tex
\section{Метасвидетельства относительно решения}\label{sec:meta-evidence}
Сертификат корректного обучения не является сертификатом экономического выигрыша. Различим стоимость выбора контроллера, стоимость получения сведений о его стоимости и достаточность этих сведений для решения. Конечная метазадача сама является экземпляром модели исследования раздела~3; её сценарии --- пары задач с фиксированными размерами носителей.

\subsection{Выбор финансируемых исследовательских планов}
План задаёт список базовых схем правил, лимит обучения $\lambda$, резерв завершения $R$ и начальную плату $k$. Сначала выполняется оплачиваемый префикс обучения в пределах $\lambda$, затем текущий фронтир, неизрасходованный лимит и $R$ передаются статическому исполнителю с отсечением. Требуется доказательство, что передаваемый остаток не меньше статического потенциала обязательств. Запись префикса, включая любое ожидающее событие обучения, возвращается отдельно; ожидающая компиляция не является оставшимся обязательством по ответу. При нулевом лимите обучение отключено, включая плату за поиск в библиотеке.

\begin{theorem}[Финансируемые планы и метавыбор]\label{thm:meta-selection}
Каждый финансируемый план возвращает в точности каноническое членство для своей пары задач. Его фактическая модельная стоимость, включая начальную плату, не превосходит $k+\lambda+R$. Пусть для конечного каталога таких планов, индексированного действиями, таблица $U(a,w)$ ограничивает сверху фактическую стоимость завершения $C(a,w)$ каждого плана в каждом сценарии, и задана фиксированная плата контроллера $k_0$. Оптимизатор исследования раздела~3 минимизирует $k_0+W_E(t)$ при терминальной потере $U$; его прогнозируемая худшая стоимость ограничивает фактическое исполнение. Если всюду $U=C$, оптимизатор минимизирует фактическую худшую стоимость по допустимым метадеревьям.
\end{theorem}
\begin{proof}
Префикс обучения сохраняет корректность. Финансируемая передача завершается по теореме~\ref{thm:funded-search}; сохранение средств в обеих фазах даёт оценку и отделяет неизрасходованный резерв от затрат. На метауровне выбранное дерево исполняет только выбранный им план. В каждом мире его фактическая стоимость равна сумме реализованных плат за запросы и $C(a,w)$, ограниченной соответствующей суммой с $U$. Теорема~\ref{thm:investigation} оптимизирует эту верхнюю целевую функцию; точность таблицы превращает неравенство в фактическую оптимальность. См. \rc{perform\_complete}, \rc{actualBill\_le\_reservation}, \rc{optimizer\_bound\_optimal} и \rc{optimizer\_actual\_optimal} в \rc{MetaInvestigation.lean}.
\end{proof}
Таблица, доказательства финансирования и плата контроллера предоставляются, а не выводятся из бесплатного оракула. Явный свидетель уменьшает худшую модельную стоимость с $53$ до $25$ (включая объявленную плату метанаблюдения); офлайн-получение полной таблицы стоит $465$. При этих фиксированных платах строгая окупаемость требует не менее $17$ использований, поскольку $465+25n<53n$ равносильно $465<28n$. Это модельный свидетель, не прогноз будущей нагрузки. Необходимость получить таблицу стоимости мотивирует частичные свидетельства.

\subsection{Интервальные свидетельства и платное уточнение}
При фиксированных мирах, портах, ценах и глубине каждая неизвестная терминальная стоимость $C(a,w)$ представляется интервалом, индексированным доказательством:
\[
  L(a,w)\le C(a,w)\le U(a,w).
\]
Интервалы можно пересекать; нижние границы повышаются, верхние понижаются, фактическая стоимость остаётся покрытой. Пусть $V_L,V_C,V_U$ --- соответствующие минимаксные значения раздела~3, $W_C$ --- худшая стоимость при фактических потерях, а $t_U$ --- оптимизатор верхней модели.
\begin{theorem}[Достаточный интервальный сертификат]\label{thm:interval}
$V_L\le V_C\le V_U$, и для каждого допустимого конкурирующего дерева $t$
\[
 W_C(t_U)\le W_C(t)+(V_U-V_L).
\]
Зазор $V_U-V_L$ не может увеличиваться при уточнении интервалов. В частности, зазор не более $\epsilon$ достаточен для сертификации верхнего оптимизатора как $\epsilon$-близкого к оптимальному относительно фактической минимаксной цели.
\end{theorem}
\begin{proof}
Стоимости деревьев монотонны по терминальным потерям, а допустимость не меняется; поэтому соответствующие оптимальные значения упорядочены. Далее $W_C(t_U)\le V_U\le W_C(t)+V_U-V_L$, поскольку $V_L\le V_C\le W_C(t)$. Уточнение улучшает обе границы в соответствующих направлениях. Формальные утверждения: \rc{lower\_le\_true}, \rc{true\_le\_upper}, \rc{certified\_gap} и \rc{gap\_monotone} в \rc{MetaEvidenceAcquisition.lean}.
\end{proof}
Получение и оценка --- отдельные операции с положительными платами в предоставленном конечном расписании. Граница, индексированная фактической стоимостью, приобретается только после оплаты её котировки; исполнитель регистрирует платежи, монотонно обновляет свидетельства и останавливается при достаточном сертификате или исчерпании бюджета/расписания. Теорема \rc{assess\_complete} устанавливает полноту теста зазора, а не необходимое условие близости к оптимуму. Для конкретного семейства финансируемых планов приобретение шести из восьми точных ячеек сертифицирует значение $22$ при предоплате $2309$ против $3051$ за полное получение. Это свидетель расписания, установленного постфактум, а не оптимальный алгоритм приобретения. Предоплата по резерву не тождественна реализованной стоимости эксперимента.

\subsection{Достаточность для решения и оптимальное приобретение свидетельств}
Неопределённость стоимости не влечёт неопределённости решения. Например, в двух возможных моделях со стоимостями $(0,1)$ и $(100,101)$ для двух контроллеров первый оптимален в обеих, несмотря на большой разброс оптимальных значений. Поэтому используем конечный непустой каталог $\mathcal C$ допущенных контроллеров и список $E$ совместимых гипотез. Каждая гипотеза $h$ задаёт стоимость $J_h(c)$ и эталон $v_h\le J_h(c)$ для всех $c\in\mathcal C$. Эталон может быть достигаемым минимумом каталога, но не обязан быть оптимумом по более широкому классу. Определим
\[
 r_h(c)=J_h(c)-v_h,\qquad R(E)=\min_{c\in\mathcal C}\max_{h\in E}r_h(c),
\]
с максимумом пустого списка, равным нулю.
\begin{theorem}[Достаточность относительно решения]\label{thm:decision}
$R(E)\le\epsilon$ тогда и только тогда, когда некоторый допущенный контроллер удовлетворяет $J_h(c)\le v_h+\epsilon$ для каждого $h\in E$. Для эталона каталога это эквивалентно $J_h(c)\le J_h(c')+\epsilon$ для всех $h\in E,c'\in\mathcal C$. Уточнение списка гипотез не увеличивает $R$.
\end{theorem}
\begin{proof}
Минимум достигается конечным выбором, а ограничение максимума эквивалентно поточечным неравенствам. Для каталога в одну сторону используется нижняя граница минимума каталога на всех конкурентах, в обратную --- достигающий минимум конкурент для каждой гипотезы. Уточнение подмножеством может лишь уменьшать худшее сожаление каждого контроллера, а значит, и минимум. Это \rc{residual\_iff}, \rc{catalogue\_robust\_iff} и \rc{residual\_subset} в \rc{DecisionEvidence.lean}.
\end{proof}
Попарные границы стоимости могут сертифицировать решение без идентификации модели стоимости. Принадлежность действительной гипотезы --- явная внешняя предпосылка: детерминированный ответ порта $q$ фильтрует $E$, сохраняя действительную гипотезу, если она присутствовала изначально (\rc{refine\_truth}). Используя терминальную потерю $r_h(c)$ в модели исследования, получаем оптимальное дерево приобретения по сумме затрат на запросы и сожаления в худшем случае, с терминальным значением $R(E)$ и рекурсией Беллмана раздела~3. Это \rc{acquisition\_optimal} и \rc{acquisition\_bellman}. Фактическая полная стоимость при внешней подготовке $k$ в точности равна $k+v_h+$ реализованная сумма платы за запросы и сожаления; минимизация худшего сожаления в общем случае не минимизирует худшую фактическую стоимость, если $v_h$ меняется (\rc{total\_bill\_identity}).

Жёсткий контракт отличается от мягкой суммы. Используя ограниченную конструкцию теоремы~\ref{thm:investigation}, \rc{hardAssess} успешен тогда и только тогда, когда существует допустимое дерево приобретения с расходами на наблюдения не более $b$ и терминальным сожалением не более $\epsilon$ для каждой совместимой гипотезы. При достижимости его дерево также минимизирует худшие расходы на наблюдения среди допустимых деревьев, удовлетворяющих допуску сожаления (\rc{hard\_assess\_complete}, \rc{hard\_minimal\_observation}). Эти гарантии требуют принадлежности действительной гипотезы и заданного конечного каталога, а не автоматической генерации гипотез или бесплатной валидации. В частности, предварительный выбор контроллера не разрешает его установку в частично исследованный текущий фронтир. Это отдельная задача замены раздела~6.
% END ORIGINAL: source/ru/section5.tex
\clearpage
% BEGIN ORIGINAL: source/ru/section6.tex
\section{Сертифицированная замена контроллера}\label{sec:replacement}
Предыдущие конструкции отделяют существование переносимого исследования от стоимости его обнаружения. Однако выбор контроллера до начала поиска не обосновывает его замену в ходе частично завершённого поиска. Замена должна сохранять и неразрешённые обязательства, и уже сертифицированные результаты; её оплата не должна лишать возможности завершить исследование. Зафиксируем индексированные исходную и целевую задачи $S,T$ и канонический эталонный список $\CanonRef=\rc{discoverAll}(S,T)$ раздела~3. В разделах~6--7 учётные величины означают натуральные модельные платы, а не измеренное время или память.

\subsection{Продолжения и их денотация}
Контроллер $C=(\mathcal R,\tau)$ состоит из библиотеки сертифицированных правил отсечения и базовых тарифов поиска. Фронтир $f=(J,A,\Lambda)$ содержит ожидающие задания $J$, список сертифицированных морфизмов $A$ и список поисковых квитанций $\Lambda$. Обозначая через $\Cov(x,J)$ покрытие морфизма ожидающим заданием, определим корректность фронтира в точности как
\begin{equation}\label{eq:frontier-correct}
 \Correct(f)\;\Longleftrightarrow\;
 (A\subseteq\CanonRef)\ \land\
 \bigl(\forall x\in\CanonRef,\ x\in A\ \lor\ \Cov(x,J)\bigr),
\end{equation}
где включение и членство относятся к спискам и не утверждают уникальность или непересечение. Таким образом, фронтир представляет и установленное знание, и оставшиеся обязательства поиска. Для фронтиров $f,g$ определим эквивалентность продолжений:
\begin{equation}\label{eq:same-cont}
 f\simeq g\;\Longleftrightarrow\;
 A_g=A_f\ \land\ \Lambda_g=\Lambda_f\ \land\
 \forall x:\mathrm{Morphism}(S,T),\ \Cov(x,J_g)\leftrightarrow\Cov(x,J_f).
\end{equation}
Продолжение $c=(C,f,R;\pi)$ содержит доказательство $\pi$ корректности и границы потенциала $P(C,f)\le R$. Резерв $R$ --- достаточная верхняя граница ресурсов завершения, а не расход. Определим $b(c)=\bill(c)$ как фактическую модельную стоимость исчерпывающего исполнения текущего статического продолжения из остатка $R$ и нулевых дополнительных расходов. По финансируемому завершению $b(c)\le R$. Действующее состояние также требует $R\le B$, где $B$ --- доступный остаток средств.

\begin{theorem}[Замена с сохранением продолжения]\label{thm:rebase}
Замена $\mathcal R$ другой сертифицированной библиотекой $\mathcal R'$ при том же $\tau$, переинициализация ожидающих очередей правил и предоставление корректного нового резерва дают продолжение $c'$ с $f\simeq f'$. Его завершение имеет в точности то же каноническое членство, что и завершение $c$; членство обоих равно членству в $\CanonRef$.
\end{theorem}
\begin{proof}
Операция переинициализации сбрасывает только очередь правил ожидающего задания области портов. Она сохраняет его область, каждое задание действий и его необработанный суффикс, точно копирует $A,\Lambda$. Покрытие не зависит от очереди правил, поэтому инвариантно покрытие каждого задания и, следовательно, их списка. Отсюда \eqref{eq:same-cont}, а \eqref{eq:frontier-correct} переносится на новый фронтир. Финансируемое завершение каждого продолжения даёт членство в $\CanonRef$ независимо от выбранной библиотеки правил. Это \rc{rebase\_same}, \rc{same\_correct} и \rc{completion\_full} в \rc{ControllerReplacement.lean}.
\end{proof}
Эта эквивалентность сильнее равенства множеств конечных ответов: она сохраняет денотацию оставшихся обязательств без перезапуска поиска. Она не разрешает произвольную миграцию исполнимого кода или перенос полного состояния обучения.

\subsection{Ликвидность и сравнительные свидетельства стоимости}
Абсолютное предложение индексировано точным текущим продолжением $c$ и содержит новый контроллер, доказательства неизменности тарифов и финансирования, плату за миграцию $K$ и доказанные границы $L\le b(c)$ и $b(c')\le U$. Пусть $H$ --- плата за проверку. Числовой механизм допуска принимает предложение тогда и только тогда, когда
\begin{equation}\label{eq:absolute-gate}
 H+R\le B,\qquad H+K+R'\le B,\qquad H+K+U<L.
\end{equation}
Первое неравенство защищает завершение текущим контроллером даже при неудаче проверки. Второе защищает завершение после миграции. Последнее устанавливает строгий выигрыш с учётом платы за проверку и миграцию. Низкая оставшаяся стоимость не означает доступности замены, поскольку резерв отличен от стоимости.

Независимые границы могут терять существенные корреляции. Поэтому конструкция относительных свидетельств заменяет одну пару $(L,U)$ конечным списком $E$ совместных пар нижних и верхних границ с предоставленным доказательством покрытия действительной пары:
\begin{equation}\label{eq:actual-coverage}
 \exists (\ell,u)\in E:\quad \ell\le b(c)\quad\land\quad b(c')\le u.
\end{equation}
Это совместный сертификат, а не список несвязанных маргинальных границ. В частности, $E$ не может быть пустым на типизированной границе, хотя необработанный универсальный арифметический тест принял бы пустой список. Предикат стоимости имеет вид
\begin{equation}\label{eq:relative-gate}
 \operatorname{relativeOK}(H+K,E)\;\Longleftrightarrow\;
 \forall(\ell,u)\in E,\ H+K+u<\ell.
\end{equation}
Для одноэлементного списка $[(L,U)]$ восстанавливается абсолютный экономический тест \eqref{eq:absolute-gate}. Например, пары $(7,1)$ и $(70,64)$ проходят при $H+K=2$, тогда как независимые экстремумы $(7,64)$ не проходят. Вторая строка свидетеля релиза представляет синтетическую неопределённость, а не измеренный второй мир.

\begin{theorem}[Точный допуск и корректность выигрыша]\label{thm:gate}
Для типизированного относительного предложения фиксированный числовой механизм допуска возвращает \emph{accepted} (принято) тогда и только тогда, когда выполнены первые два неравенства \eqref{eq:absolute-gate} и \eqref{eq:relative-gate}. При принятии $H+K+b(c')<b(c)$. Применение сохраняет эквивалентность продолжения, базовые тарифы поиска и финансируемый резерв.
\end{theorem}
\begin{proof}
Эквивалентность следует из последовательных ветвей \rc{gate}; универсальный тест строк отражён в булево значение посредством \rc{relativeOK\_iff}. Выберем покрывающую пару из \eqref{eq:actual-coverage}. Тогда $H+K+b(c')\le H+K+u<\ell\le b(c)$, что доказывает строгий выигрыш. Переинициализация использует теорему~\ref{thm:rebase}. Остаток после принятия $B-H-K$ не меньше $R'$ по второму условию. При отклонении сохраняется прежнее продолжение, а первое условие ограничивает его резерв остатком после проверки. При отсрочке состояние не меняется. Это \rc{gate\_accepted\_iff}, \rc{relative\_sound}, \rc{apply\_same} и \rc{apply\_liquid}.
\end{proof}

Публичная попытка использует $H=h+1$ для базовой платы $h\in\Nat$. Её вердикт --- обычный четырёхзначный результат; применение требует доказательства, что этот вердикт является выходом фиксированного механизма допуска для данного точного состояния и предложения. Таблица~\ref{tab:debit} задаёт операционный учёт. Первое нарушенное условие определяет вердикт; тест строк вычисляется только после прохождения обоих условий резерва.
\begin{table}[htbp]
\centering\small
\caption{Учёт допуска. $d$ --- списание, а не резерв завершения.}\label{tab:debit}
\begin{tabularx}{\linewidth}{@{}lXll@{}}\toprule
Вердикт & Условие (по порядку) & Установлено & $d$\\\midrule
Отсрочка & Не выполнено условие прежнего резерва & $c$ & $0$\\
Отказ: ликвидность & Не выполнено условие нового резерва & $c$ & $H$\\
Отказ: выигрыш & Не пройден тест совместных пар & $c$ & $H$\\
Принято & Пройдены все три теста & $c'$ & $H+K$\\\bottomrule
\end{tabularx}
\end{table}

Пусть сессия $s$ содержит действующее состояние и относительный журнал. Обозначим $D(s)=\spent(s)$ как сумму поисковых квитанций, плат за прежние попытки и плат относительного журнала, а $W(s)=B(s)+D(s)$ как общие средства. Применение вычитает ровно $d$ из остатка и добавляет $d$ в журнал, сохраняя $W$. Предложение уже несёт доказательства; механизм допуска не оценивает неизвестные стоимости бесплатными пробными запусками и не ищет доказательств. Он проверяет объявленные числовые условия. Внешние платежи за свидетельства регистрируют расход, но сами по себе не удостоверяют истинность или приобретение свидетельства.

\subsection{Закон смешанных историй замен}
Для финансируемой сессии определим целевую функцию завершения
\begin{equation}\label{eq:root-objective}
 \Phi(s)=D(s)+b(c_s).
\end{equation}
Это полные модельные расходы при завершении текущего продолжения, а не онлайн-оценка. Разрешённые действия --- финансируемый префикс текущего поиска, попытка, индексированная текущим продолжением, и доступная по цене внешняя квитанция о свидетельстве. Связная трасса $s_0\to\cdots\to s_n$ состоит из фактических применений действий с равенством соседних концов; это не список произвольно заявленных постусловий. В частности, разрешимое равенство произвольных состояний, содержащих функции, не предполагается.

Для каждого шага определим чистый выигрыш $g_i$ и накладные расходы $o_i$. Принятая замена имеет $g_i=\Phi(s_{i-1})-\Phi(s_i)>0$ уже за вычетом проверки и миграции, и $o_i=0$. Оплаченный отказ или плата за свидетельство имеет $g_i=0$, а $o_i$ равно соответствующему платежу. У префикса поиска и отсрочки обе величины нулевые. Пусть $G=\sum_i g_i$, $O=\sum_i o_i$, $N_p$ --- число неотложенных попыток замены, $N_a$ --- число принятых.

\begin{theorem}[Композиционный учёт корневых замен]\label{thm:root-trace}
Для каждой связной конечной корневой трассы
\begin{align}
 \Phi(s_n)+G&=\Phi(s_0)+O,\label{eq:trace-objective}\\
 W(s_n)&=W(s_0),& R(s_n)&\le B(s_n),\label{eq:trace-funds}\\
 N_p+B(s_n)&\le B(s_0),& N_a&\le G.
\end{align}
Кроме того, отношение концов $\Pres(s_0,s_n)$ точно сохраняет базовые тарифы поиска и прежнюю историю попыток, список сертификатов --- как префикс, а поисковые квитанции и относительный журнал --- как суффиксы списков в обратном хронологическом порядке.
\end{theorem}
\begin{proof}
Для префикса фиксированного контроллера детерминированное продолжение даёт $b(c)=d+b(c')$, где $d$ --- затраты префикса. Поэтому \eqref{eq:root-objective} не меняется. Отложенная попытка тождественна. Отклонённая попытка или квитанция свидетельства сохраняет $c$ и увеличивает $D$ на записанный платёж. Принятие строго уменьшает $\Phi$ по теореме~\ref{thm:gate}; эта разность определяет чистый выигрыш. Следовательно, каждый шаг удовлетворяет $\Phi(s_i)+g_i=\Phi(s_{i-1})+o_i$, сохранению средств и указанному отношению истории. Отношения префикса сертификатов и суффикса квитанций компонуются конкатенацией списков. Индукция по связной трассе складывает обе неотрицательные статистики на каждом шаге; положительные тарифы проверки ограничивают число оплаченных попыток, каждый принятый выигрыш не меньше единицы. Тождества: \rc{remainingBill\_prefix}, \rc{action\_laws}, \rc{laws\_compose} и \rc{trace\_laws} в \rc{RecursiveTrace.lean}.
\end{proof}
Поскольку $D(\finish(s))=\Phi(s)$, \eqref{eq:trace-objective} даёт необходимое и достаточное условие строгой чистой экономии относительно немедленного завершения текущим контроллером:
\begin{equation}\label{eq:net-gain}
 D(\finish(s_n))<D(\finish(s_0))\quad\Longleftrightarrow\quad O<G.
\end{equation}
Это \rc{trace\_final\_gain\_iff}. Теорема допускает убытки от исследовательских попыток без ослабления корректности. Она выявляет следующее препятствие: сам выбор момента таких попыток может быть неадекватен. Это мотивирует явное представление политики предложений как заменяемого объекта вместо фиксированной последовательности эвристических вмешательств.
% END ORIGINAL: source/ru/section6.tex
% BEGIN ORIGINAL: source/ru/section7.tex
\section{Композиционное замыкание метаполитик}\label{sec:closure}
\subsection{Оплачиваемая и заменяемая политика}
Теперь отличим корневой контроллер от конечной политики, решающей, следует ли пытаться его заменить. Зафиксируем текущее продолжение $c_*$ и поставщика $p$, индексированного $S,T,c_*$. Поставщик содержит детерминированное отложенное булево наблюдение, отложенное вычисление типизированного относительного предложения для $c_*$ и объявленные базовые платы за наблюдение, подготовку и проверку. Фрейм $f$ содержит сессию $s$ с доказательством $c_s=c_*$. Политика $t$ --- конечное дерево с листьями \emph{keep} (сохранить) или \emph{attempt} (попытаться заменить) и бинарными узлами наблюдения. Конкретный язык использует один порт наблюдения и разрешает повторные зондирования; он не обязан принадлежать классу допустимых деревьев с однократным использованием портов из прежней задачи оптимального синтеза.

Интерпретация осуществляется структурной рекурсией. Зондирование и подготовка вызываются только в финансируемых ветвях после списания соответствующей платы $h+1$ с сохранением корневого резерва. Терминальная попытка подготавливает предложение и вызывает фиксированный относительный механизм допуска раздела~\ref{sec:replacement}; невыполнение условия блокирует соответствующую операцию. Интерпретация регистрирует фактические булевы ответы, подготовленные типизированные предложения и вердикты. Она выдаёт сессию $I_p(t,f)$ с сохранёнными историей и средствами независимо от принятия замены. Затем политика завершается финансируемым корневым исполнителем. Определим её полную стоимость (включая уже понесённые расходы):
\begin{equation}\label{eq:policy-bill}
 \PB_p(t,f)=\Phi(I_p(t,f))=D(\finish(I_p(t,f))).
\end{equation}
Соответствующие определения и законы содержатся в \rc{MetaPolicy.lean}. В частности, \rc{policyBill\_complete} доказывает второе равенство.

Машина $m=(f,t,L)$ добавляет к фрейму текущую политику и её историю оплаченных обновлений $L$. Рассмотрим предложенную замену $t'$ с базовой платой подготовки $a$ и $q=a+1$. Перед вызовом фабрики предложения операция обновления проверяет $q+R\le B$. Если условие не выполнено, возвращается та же машина, без платы и вызова. Иначе списывается $q$, получается новый фрейм $f^+$ с прежним продолжением и вызывается фабрика. Предоставленное предложение доказывает
\begin{equation}\label{eq:policy-bounds}
 \lambda\le\PB_p(t,f),\qquad \PB_p(t',f^+)\le\upsilon.
\end{equation}
Фиксированный механизм допуска устанавливает $t'$ в точности при $\upsilon<\lambda$, иначе сохраняет $t$ во фрейме $f^+$. Оба оплаченных исхода записывают предшествующее и предложенное дерево, границы и плату. Плата подготовки уже учтена фреймом после платежа в \eqref{eq:policy-bounds}; её повторное добавление означало бы двойной учёт. Поставщик и корневой контракт не меняются.

\begin{theorem}[Локальное улучшение политики]\label{thm:policy-local}
Каждое выполненное обновление сохраняет $W$, $\Pres$, ликвидность и предшествующую историю обновлений. Принятое обновление из $m$ в $m'$ строго уменьшает $\PB_p(t,f)$ относительно точного значения до обновления.
\end{theorem}
\begin{proof}
Платёж --- доступная по средствам плата за свидетельство, поэтому сохраняет продолжение и средства. Установка дерева меняет только активную политику; не исполняет её и не изменяет текущего продолжения. При принятии цепочка $\PB_p(t',f^+)\le\upsilon<\lambda\le\PB_p(t,f)$ даёт строгое уменьшение. Отсроченное обновление тождественно; отклонённое сохраняет активное дерево, но не прежний остаток. См. \rc{upgrade\_funds}, \rc{upgrade\_preserves}, \rc{upgrade\_history} и \rc{upgrade\_gain}.
\end{proof}
Сравнение использует сертифицированные границы точных фреймов до и после платежа, а не предполагаемый бесплатный оптимизатор. Фабрики и их доказательные обязательства являются входами. Модельная плата подготовки не есть эмпирическая оценка вычислительных усилий по построению произвольных доказательств.

\subsection{Замыкание конечных историй обновлений}
План содержит базовую плату и фабрику, производящую предложение с доказательствами для фактической текущей машины при наличии средств. Переход записывает предсостояние, план и действительный результат обновления с равенством, удостоверяющим исполнение. Допустимая трасса из $n$ переходов соединяет $m_0$ с $m_n$ через эти результаты. Пусть $Q_\Sigma$ --- сумма списаний за подготовку (нуль при отсрочке, иначе $a+1$), а $N_p,N_a$ --- числа оплаченных и принятых обновлений трассы. Обозначим $\PB_p(t_i,f_i)$ через $\Psi_i$.

\begin{theorem}[Конечное замыкание метаполитик]\label{thm:meta-closure}
Каждая допустимая конечная трасса с фиксированным $p$ удовлетворяет
\begin{align}
 D(s_n)&=D(s_0)+Q_\Sigma,& W(s_n)&=W(s_0),\\
 R(s_n)&\le B(s_n),& N_a&\le N_p\le Q_\Sigma,\\
 N_p+B(s_n)&\le B(s_0).
\end{align}
Она сохраняет $\Pres(s_0,s_n)$, а конечный список оплаченных обновлений расширяет начальный префиксом длины $N_p$. Если ни одно обновление в трассе не имеет вердикта \emph{rejected} (отклонено), то
\begin{equation}\label{eq:meta-gain}
 \Psi_n+N_a\le\Psi_0.
\end{equation}
\end{theorem}
\begin{proof}
Локальный закон учёта добавляет ровно списание за подготовку и одну оплачиваемую запись истории в каждой финансируемой ветви. Индукция по трассе компонует эти тождества и транзитивное отношение сохранения; сохранение средств даёт границу остатка. В трассе без отказов каждый шаг либо тождествен, либо уменьшает натуральную целевую функцию $\Psi$ не менее чем на единицу, поэтому суммирование локальных неравенств доказывает \eqref{eq:meta-gain}. Это \rc{trace\_laws}, \rc{trace\_paid\_balance}, \rc{trace\_liquid} и \rc{trace\_no\_rejection\_gain} в \rc{MetaClosure.lean}.
\end{proof}

Условие об отсутствии отказов в \eqref{eq:meta-gain} существенно. Корневое тождество \eqref{eq:trace-objective} нельзя повторно использовать как закон только плат для $\Psi$: отклонённое метаобновление уменьшает ликвидность и может изменить набор операций, доступных впоследствии интерпретируемому дереву. В замкнутом свидетеле препятствия по глубине текущее продолжение стоит 7, замена стоит 1, резерв равен 100, подготовка, проверка и миграция стоят по 1. При остатке 103 попытка завершается с полной затратой 4. Оплата одной единицы за неудачное метаобновление оставляет 102. Подготовка и корневая проверка ещё доступны, миграция --- нет; корневой допуск отказывает по ликвидности и завершение прежнего продолжения даёт полные расходы $1+1+1+7=10$, а не $4+1$. Это отрицательный контроль \rc{Release.Witnesses.low\_liquidity}. Он устанавливает разрывность, обусловленную условием финансирования, а не нарушение сохранения ресурсов.

\subsection{Исполнимая теорема замыкания}
Релиз принимает конечный список типизированных планов (или предоставленный связный запуск объявленной длины), исполняет обновления, однократно интерпретирует последнее установленное дерево и завершает полученную корневую сессию за счёт защищённого резерва. Запишем
\begin{equation}\label{eq:release-pipeline}
 s_f=\finish\bigl(I_p(t_n,f_n)\bigr),\qquad
 m_n=\operatorname{runPlans}(p,m_0,\mathcal L).
\end{equation}
Здесь $m_n$ означает конечную машину возвращённого запуска, $n=|\mathcal L|$, а полный результат сохраняет запуск, интерпретированный эпизод и конечную сессию со связывающими равенствами. Модель расписания намеренно конечна: каждый план расходует одну ячейку списка, даже если условие финансирования откладывает его. Интерпретация дерева завершается структурно, а корневое завершение --- по своей финансируемой границе. Типизированный интерфейс требует изначально финансируемого продолжения; он не классифицирует нефинансируемые входы как невозможные.

\begin{theorem}[Финансируемый релиз с доказательствами]\label{thm:release}
Для фиксированных индексированных $S,T$, текущего продолжения $c_*$, поставщика $p$, допустимой начальной машины $m_0$ и конечного корректно типизированного списка планов $\mathcal L$ релиз строит замкнутый запуск длины $n=|\mathcal L|$, удовлетворяющий
\begin{align}
 &J_f=[] ,\qquad \forall x:\mathrm{Morphism}(S,T),\quad x\in A_f\iff x\in\CanonRef,\label{eq:release-complete}\\
 &A_f=[]\iff \neg\operatorname{Nonempty}(\mathrm{Morphism}(S,T)),\label{eq:release-impossible}\\
 &W(s_f)=W(s_0),\qquad \Pres(s_0,s_f),\qquad D(s_f)=\Psi_n.\label{eq:release-conserve}
\end{align}
В \eqref{eq:release-impossible} правая часть означает необитаемость типа морфизмов. Если в трассе обновлений нет отклонённых обновлений, то $D(s_f)+N_a\le\Psi_0$.
\end{theorem}
\begin{proof}
Структурная рекурсия по списку строит связную трассу и доказывает её длину. По теореме~\ref{thm:meta-closure} полученная машина сохраняет средства и историю. Интерпретация снова их сохраняет, а финансируемое завершение исполняет все ожидающие задания. Теорема полноты финансируемого поиска даёт \eqref{eq:release-complete}. Если морфизм существует, конструктивное обнаружение даёт канонического представителя с теми же картами портов и действий; он обязан входить в $A_f$. Обратно, каждый элемент списка населяет тип морфизмов. Это доказывает \eqref{eq:release-impossible}. Композиция трёх законов сохранения и \eqref{eq:policy-bill} даёт \eqref{eq:release-conserve}. Последнее неравенство следует из \eqref{eq:meta-gain}. Соответствующие цели проверки ядром: \rc{result\_steps}, \rc{result\_exhausted}, \rc{result\_complete}, \rc{result\_impossible\_iff}, \rc{result\_funds}, \rc{result\_history}, \rc{result\_bill} в \rc{Release.lean}, и \rc{closed\_gain} в \rc{MetaClosure.lean}.
\end{proof}

Утверждение о завершении касается точного членства канонических морфизмов, а не перечисления каждой возможной карты миров. Напечатанная сводка --- проекция типизированного результата, а не декодер сертификатов или утверждение о долговременном хранении. Теорема замыкания также не влечёт автономного синтеза любой фабрики предложения, произвольного самопереписывания или неограниченного улучшения. Она устанавливает ограниченную композицию допущенных изменений при фиксированном контракте с полным корневым завершением, даже когда предложенные улучшения не проходят. Следующий раздел проверяет релиз на явных положительных и отрицательных вычислительных свидетелях, различая проверку ядром и исполнение модели.
% END ORIGINAL: source/ru/section7.tex
\clearpage
% BEGIN ORIGINAL: source/ru/section8.tex
\section{Формальная верификация и исполнимые свидетели}\label{sec:verification}
Утверждения разделов~\ref{sec:investigation}--\ref{sec:closure} опираются на три различных вида свидетельств: термы доказательств, проверенные ядром Lean, исполнимые примеры конечной модели и независимые эталонные вычисления. Ни одно из них не является измерением физического быстродействия, а согласие конечного числа запусков не заменяет общих утверждений, доказанных в формализации.

\subsection{Механизация и границы свидетельств}
Принятый канонический стек использует Lean~4.36.0 \cite{lean4}, закреплённый за коммитом\\
\texttt{611c117bfaf7b7ae1cfcf246ce6fba31e8f509e9}. Он включает 21 исходный файл Lean (включая CLI) и 542 явно аудируемые декларации: 520 теорем и 22 определения. Таблица~\ref{tab:formal-audit} разбивает цели аудита по группам зависимостей архитектуры; эти числа не являются числом независимых результатов или мерой сложности доказательств. Последние два файла предоставляют типизированный релиз и его исполнимую проекцию. Две ранние исторические ветви сохранены отдельно и не импортируются этим стеком.

\begin{table}[htbp]
\centering\small
\caption{Соответствие артефактов: исторические идентификаторы этапов, математические конструкции и цели аудита. Раздел~2 имеет отдельные исторические доказательные артефакты.}\label{tab:formal-audit}
\begin{tabularx}{\linewidth}{@{}lXrr@{}}\toprule
Этап & Основная конструкция & Файлы & Цели\\\midrule
T19--T22 & Исследование, перенос, библиотека, открытие & 4 & 132\\
T22.9--T24.7 & Платный поиск, отсечение, продолжение, ядра, обучение & 6 & 119\\
T25--T27 & Метаисследование и свидетельства для решения & 3 & 101\\
T28--T29c & Замена контроллера и законы трасс & 4 & 132\\
T30--T31 & Интерпретация политик и замыкание & 2 & 40\\
T32 & Релиз и CLI & 2 & 18\\\midrule
Всего & & 21 & 542\\\bottomrule
\end{tabularx}
\end{table}

Канонический исполнитель проверяет полный манифест и 19 закреплённых унаследованных исходников, компилирует в чистом временном каталоге импортов, аудирует ожидаемые полные имена и зависимости каждой заданной декларации, запускает 13 независимых скриптов конечных моделей и сравнивает фактический вывод CLI с отдельным эталоном Python. Единственные разрешённые аксиоматические зависимости 542 целей --- \rc{propext} и \rc{Quot.sound}; предоставленный успешный журнал фиксирует ноль ошибок и предупреждений компилятора, отсутствие зависимостей от \rc{sorryAx} и \rc{Classical.choice} в этих отчётах. Два анализатора требуют точных имён отчётов и отклоняют неполные отчёты, дубликаты, неразрешённые зависимости и неуспешные коды выхода компилятора. Десять тестовых примеров анализатора аудита и двенадцать примеров анализатора проекций проверяют эти механизмы; это тесты проверяющей программы, а не дополнительные запуски ядра.

Принятый успешный журнал предоставлен автором. В среде первой сборки архива не было Lean, и запуск остановился с кодом~127. Позднее предоставленные полные журналы модулей и исполнителя побайтно совпадают с соответствующим прежним успешным журналом. Это согласуется с детерминированным повтором, но независимо не устанавливает отдельного нового процесса компиляции. Согласование сохраняет это различие в \rc{catalog/VERIFICATION.json} со статусом \rc{SUPPLIED\_LOGS\_AND\_PYTHON\_REPLAY\_PASS}. Последующий повторный аудит согласованного архива (повторённый при подготовке английской редакции RC4) проверил все 21 закреплённый исходник и 542 имени отчётов, 21 журнал модулей и девять проецируемых сценариев релиза; все 45 команд Python завершились с нулевым кодом. Этот аудит не вызывал Lean. Совпадение хешей устанавливает идентичность содержимого, а не независимое удостоверение исполнения. Историческая нумерация ограничена этой картой артефактов. Ранние исходники и первичные связи T03 и T14--T18 сохранены отдельно; для T01, T02 и T04--T13 нет однозначно установленных первичных связей с этапами. Они не являются предпосылками канонического доказательства релиза. См. \rc{GAPS\_AND\_STATUS\_RU.md}.

\clearpage
\subsection{Положительные и отрицательные вычислительные свидетели}
Девять сценариев релиза проверяют составной исполнитель, а не изолированный предикат допуска. Для каждого независимый эталон строит канонический список ответов, проверяет остатки, упорядоченные квитанции поиска и обновлений, события политики и прежние сертификаты. Сохранённая проекция Lean CLI совпадает с эталоном. Результат \emph{impossible} означает исчерпание с пустым каноническим списком в заданном классе морфизмов, а не невозможность для произвольных алгоритмов.

\begin{table}[htbp]
\centering\small
\caption{Свидетели релиза из \rc{FINITE\_RELEASE.json} и предоставленного журнала CLI. Все начинаются с нулевых записанных расходов, кроме случая Predecessor с пятью единицами расходов. $B_0$ --- начальный доступный остаток, $D_f$ --- конечные совокупные расходы, $B_f$ --- конечный остаток. Последний столбец --- число канонических сертификатов.}\label{tab:release-cases}
\begin{tabularx}{\linewidth}{@{}lXrrrr@{}}\toprule
Сценарий & Вердикты метаобновлений / корневой исход & $B_0$ & $D_f$ & $B_f$ & $|A_f|$\\\midrule
Baseline & нет; сохранить & 110 & 7 & 103 & 0\\
Single & принято; заменить & 110 & 6 & 104 & 0\\
Mixed & принято, отклонено; заменить & 110 & 7 & 103 & 0\\
Double & принято, принято; заменить & 110 & 6 & 104 & 0\\
Fallback & отложено, отложено; блокировка & 100 & 7 & 93 & 0\\
Liquidity baseline & нет; заменить & 103 & 4 & 99 & 0\\
Liquidity rejected & отклонено; отказ по ликвидности & 103 & 10 & 93 & 0\\
Nonempty & нет; сохранить & 100 & 11 & 89 & 2\\
Predecessor & нет; сохранить & 95 & 11 & 89 & 2\\\bottomrule
\end{tabularx}
\end{table}

\paragraph{Сквозной пример по слоям.}
Рассмотрим действительный свидетель препятствия по глубине: $S,T$ имеют по одному миру, два порта, одно терминальное действие, постоянные ложные чтения, нулевые терминальные потери и цены портов; $d_S=1$ и $d_T=0$. Тест морфизма раздела~3 не проходит ограничение глубины, поэтому $\CanonRef=[]$. Тем не менее фронтир без отсечения должен обработать три шага разбиения и четыре шага подготовки по единичной цене (стоимость $7$). Сертифицированное правило глубины отвергает начальную область за один оплаченный шаг (стоимость $1$). При $B=110$, $R=R'=100$ и $H=K=1$ условия раздела~6 допускают замену: $101\le110$, $102\le110$ и $1+1+1<7$. Корневая замена с завершением стоит $3$, не меняя требуемый результат.

Теперь сравним альтернативные исполнения на той же паре, а не последовательные стоимости. Исходная метаполитика сохраняет старый контроллер (стоимость $7$). Замкнутая модель приобретения сопоставляет ожидающий/исчерпанный фронтир и действия keep/attempt; выбранное ею дерево зондирует, остались ли задания, и пытается заменить контроллер при истинном ответе. Метаобновление оплачивается, его точные границы равны $7$ и $6$, оно устанавливает это дерево. Фактический релиз с одним обновлением платит по единице за обновление, зондирование, подготовку, проверку, миграцию и окончательный вызов правила: всего $6$, остаток $104$, пустой фронтир и сертифицированная невозможность. Это случай \emph{single} таблицы~\ref{tab:release-cases}; отдельные случаи \emph{nonempty} и \emph{predecessor} проверяют положительное каноническое членство и сохранение прежних сертификатов. См. \rc{ControllerReplacement.Witnesses}, \rc{MetaPolicy.Witnesses} и \rc{Release.Witnesses.costs}.

Случай двойного обновления начинается с политики избыточных зондирований (стоимость $8$), а не с исходной политики; её стоимость уменьшается $8\to7\to6$. Обязательное обучение проигрывает фиксированному контроллеру правила чтения во всех 515 финансируемых случаях стенда обучения. Смешанные корневые трассы завершаются со стоимостями 4 и 8 против 7; стенд замыкания проверяет 965 завершённых запусков и 2418 переходов. Пара по ликвидности ($4$ против $10$) опровергает безусловную монотонность на уровне политик. Это конечные модельные свидетели, а не выборки из прикладной генеральной совокупности или измерения физической производительности.

\newpage
\subsection{Артефакт и воспроизведение}\label{sec:artifact}
Приложение --- \rc{CRTT-v1.0-FULL-ARCHITECTURE-RECONCILED.zip} \cite{crttartifact} с неизменённым принятым каноническим содержимым, историческими исходниками, каталогами, полными предоставленными журналами, результатами повторного аудита и скриптами проверки. Его внешний SHA-256:
\begin{center}\small\texttt{e09580cdb8a8f900401390db394a7bde07bb328449b50b059f2a8cd3c175c79e}\end{center}
Полный манифест содержит 766 записей, охватывающих 767-файловую оболочку за исключением самого списка контрольных сумм. Из корня распакованного архива используйте новый выходной каталог для повторного аудита:
\begin{verbatim}
python3 verify_archive.py
python3 replay_evidence.py --output ../CRTT-log-replay
\end{verbatim}
Для свежей проверки ядром войдите в \rc{canonical/} с закреплённым компилятором Lean и выполните \verb|bash run_verify.sh|, сохраняя журнал вне этого каталога. Только последняя команда выполняет свежий запуск Lean. Приложение к рукописи включает \rc{verify\_fresh.sh} --- оболочку чистой распаковки, регистрирующую входной хеш, инструментарий, журнал исполнения и статус процесса. Она не заменяет канонический исполнитель. При подготовке английской редакции RC4 её локальное исполнение снова прошло предварительную проверку архива, но остановилось с кодом~127 до компиляции Lean. Свежая независимая проверка ядром и публичный постоянный идентификатор артефакта остаются задачами публикации; ни одна здесь не заявлена выполненной. Карты исходников связывают утверждения с декларациями и записанными строками аудита.
% END ORIGINAL: source/ru/section8.tex
% BEGIN ORIGINAL: source/ru/section9.tex
\section{Связанные работы, следствия и ограничения}\label{sec:related}
\paragraph{Информация и исследовательские решения.}
Сравнение экспериментов Блэквелла \cite{blackwell1953} даёт классическую систему упорядочения информации по качеству решений. CRTT фиксирует контракт решения и спрашивает, существует ли общее допустимое действие в каждом наблюдаемом слое; она не претендует на новое общее упорядочение статистических экспериментов. Точная рекурсия конечных деревьев наблюдения есть применение динамического программирования \cite{bellman1957}. Сложность построения оптимальных деревьев решений --- отдельная проблема \cite{hyafil1976}; теорема точного синтеза не утверждает полиномиального времени. Обучение запросами \cite{angluin1988} и адаптивная субмодулярность \cite{golovin2011} изучают стратегии приобретения при иных предпосылках и целях. Свидетель XOR раздела~3 показывает, почему неограниченное жадное приобретение здесь не обосновано. Пороговый результат раздела~2 также относится к ограниченной модели аренды и открытия, а не универсальному онлайн-оптимизатору.

\paragraph{Обучение на неудачах и проверка.}
Преобразование неудач в переиспользуемые ограничения имеет известные предшественники: Soar извлекает процедурное знание из разрешения тупиков \cite{soar1987}, а GRASP обучается ограничениям из конфликтов в пропозициональном поиске \cite{grasp1999}. Сертифицированная проверка доказательств SAT отделяет производство от валидации \cite{lrat2017}. CRTT специализирует извлечение препятствий для несовместимых строк морфизмов, доказывает область действия полученного правила и сохранение покрытия, учитывает объявленные платы за поиск, зондирование и компиляцию. Корректность обученного правила сама по себе не влечёт выгодной политики. Proof-Carrying Code (код, несущий доказательства) \cite{necula1997} служит общим архитектурным предшественником проверки предоставленного доказательства относительно фиксированной политики; в CRTT соответствующее обязательство включает денотацию незавершённого поиска, сравнительные свидетельства стоимости и резерв завершения.

\paragraph{Рефлексия и изменение управления.}
Процедурная рефлексия \cite{smith1984} делает вычислительные механизмы доступными дальнейшему вычислению. Рациональное метарассуждение \cite{russell1991} оценивает вычислительные действия по их вкладу в решения и стоимости. Динамическое обновление ПО \cite{stoyle2007} исследует корректность изменений работающих программ. Машины Гёделя \cite{schmidhuber2007} и формальные модели рациональной самомодификации \cite{everitt2016} рассматривают более широкие классы самоизменения. CRTT не включает эти подходы как частные случаи: проверяющий механизм, поставщик, контракт и базовые тарифы фиксированы, а изменяемые политики имеют явно конечный синтаксис. Конкретный вклад --- композиция конструктивного обнаружения переносов, платного обучения на препятствиях, замены с сохранением продолжения и финансируемого замыкания, а не приоритет рефлексии, обучения на конфликтах или проверки доказательств по отдельности.

\paragraph{Сопутствующие исследования.}
Предшествующие исследования автора по наблюдаемости и сохранённой памяти в Rule~110 \cite{strand110observer2026,strand110memory2026} и релевантной решениям информации в конечных каналах \cite{strandcontrol2026} мотивируют различение доступной информации и используемого управления. Они не являются ни импортируемыми доказательными зависимостями, ни эмпирической валидацией CRTT. Полная сопутствующая библиография сохранена в приложении.

\paragraph{Что должен предоставить производитель сертификатов.}
Входная граница принимает объекты с доказательствами, но не является универсальным конструктором сертификатов. Некоторые сертификаты вычисляются внутри стека: конечная сумма бюджетов заданий $P(C,f)$ даёт достаточный резерв и, следовательно, $0\le b(c)\le R$ при $R\ge P(C,f)$ (раздел~6.1). Для предоставленного финансируемого плана раздела~5.1 \rc{initialBounds} строит $[k,k+\lambda+R]$; \rc{exactRequest} вычисляет его модельную стоимость с предоплатой резерва и положительной надбавки, выдавая точечный интервал. Абсолютный сертификат стоимости вкладывается в одноэлементную совместную строку посредством \rc{fromAbsolute}. Эти построения конструктивны, но не обязательно дёшевы и не обязательно дают тесные границы. Замкнутый свидетель раздела~8.2 использует синтезированное дерево приобретения и точные границы, доказанные конечным вычислением. В общем же случае полезная нижняя граница старой стоимости, тесная верхняя граница новой, покрывающие совместные строки, принадлежность истинной гипотезы и фабрики, индексированные фактическим состоянием, должны предоставляться. Вычисление точных стоимостей само может требовать завершения; допуск избегает повторения этой работы, но не устраняет её. Здесь нет общей процедуры обнаружения таких экономичных свидетельств или оценки реального времени их производства. Таким образом, замыкание сертифицирует условную композицию, а не эффективный универсальный алгоритм открытия.

\paragraph{Область применимости и исследовательские следствия.}
Гарантии относятся к конечным носителям, детерминированным наблюдениям и корректным свидетелям. Модель не устанавливает эмпирическую истину и не оценивает автоматически построение доказательств в физических ресурсах. Релиз не мигрирует произвольное состояние программы, не переписывает ядро и не устанавливает сходимость неограниченных расписаний. Типизированная история происхождения в памяти не является устойчивым к сбоям хранением. Конечные свидетели не устанавливают обобщение на нейронных агентов или новые семейства задач. Исчерпывающее сопоставление новизны и теорема асимптотического ускорения не заявляются.

Эпистемическая интерпретация состоит в том, что результат исследования может стать проверенным оператором последующего исследования без смешения знания о корректности метода со знанием о его полезности. Отрицательные сертификаты, частичные свидетельства стоимости и эквивалентность продолжений делают явными разные стороны этого процесса. CRTT демонстрирует в конкретной исполнимой модели, как такие операторы и их управляющие политики могут пересматриваться при инвариантном обязательстве завершить то же корневое исследование. Распространение контракта на эмпирические области --- отдельная исследовательская задача, а не следствие теоремы релиза.
% END ORIGINAL: source/ru/section9.tex
\clearpage
% BEGIN ORIGINAL: source/ru/references.tex
\begin{thebibliography}{99}
\small\raggedright
\bibitem{blackwell1953} D. Blackwell. Equivalent comparisons of experiments. \emph{Annals of Mathematical Statistics}, 24(2):265--272, 1953. \href{https://doi.org/10.1214/aoms/1177729032}{doi:10.1214/aoms/1177729032}.
\bibitem{bellman1957} R. Bellman. \emph{Dynamic Programming}. Princeton University Press, 1957.
\bibitem{hyafil1976} L. Hyafil and R. L. Rivest. Constructing optimal binary decision trees is NP-complete. \emph{Information Processing Letters}, 5(1):15--17, 1976. \href{https://doi.org/10.1016/0020-0190(76)90095-8}{doi:10.1016/0020-0190(76)90095-8}.
\bibitem{angluin1988} D. Angluin. Queries and concept learning. \emph{Machine Learning}, 2:319--342, 1988. \href{https://doi.org/10.1007/BF00116828}{doi:10.1007/BF00116828}.
\bibitem{golovin2011} D. Golovin and A. Krause. Adaptive submodularity: theory and applications in active learning and stochastic optimization. \emph{Journal of Artificial Intelligence Research}, 42:427--486, 2011. \href{https://doi.org/10.1613/jair.3278}{doi:10.1613/jair.3278}.
\bibitem{soar1987} J. E. Laird, A. Newell, and P. S. Rosenbloom. SOAR: an architecture for general intelligence. \emph{Artificial Intelligence}, 33(1):1--64, 1987. \href{https://doi.org/10.1016/0004-3702(87)90050-6}{doi:10.1016/0004-3702(87)90050-6}.
\bibitem{grasp1999} J. P. Marques-Silva and K. A. Sakallah. GRASP: a search algorithm for propositional satisfiability. \emph{IEEE Transactions on Computers}, 48(5):506--521, 1999. \href{https://doi.org/10.1109/12.769433}{doi:10.1109/12.769433}.
\bibitem{lrat2017} L. Cruz-Filipe, M. Heule, W. Hunt, M. Kaufmann, and P. Schneider-Kamp. Efficient certified RAT verification. \href{https://arxiv.org/abs/1612.02353}{\textnormal{arXiv:1612.02353}}, 2016.
\bibitem{necula1997} G. C. Necula. Proof-carrying code. In \emph{Proceedings of POPL}, pp. 106--119, 1997. \href{https://doi.org/10.1145/263699.263712}{doi:10.1145/263699.263712}.
\bibitem{smith1984} B. C. Smith. Reflection and semantics in Lisp. In \emph{Proceedings of POPL}, pp. 23--35, 1984. \href{https://doi.org/10.1145/800017.800513}{doi:10.1145/800017.800513}.
\bibitem{russell1991} S. Russell and E. Wefald. Principles of metareasoning. \emph{Artificial Intelligence}, 49(1--3):361--395, 1991. \href{https://doi.org/10.1016/0004-3702(91)90015-C}{doi:10.1016/0004-3702(91)90015-C}.
\bibitem{stoyle2007} G. Stoyle, M. Hicks, G. Bierman, P. Sewell, and I. Neamtiu. Mutatis Mutandis: safe and predictable dynamic software updating. \emph{ACM Transactions on Programming Languages and Systems}, 29(4), Article 22, 2007. \href{https://doi.org/10.1145/1255450.1255455}{doi:10.1145/1255450.1255455}.
\bibitem{schmidhuber2007} J. Schmidhuber. G\"odel machines: fully self-referential optimal universal self-improvers. In \emph{Artificial General Intelligence}, pp. 199--226, Springer, 2007. \href{https://doi.org/10.1007/978-3-540-68677-4_7}{doi:10.1007/978-3-540-68677-4\_7}.
\bibitem{everitt2016} T. Everitt, D. Filan, M. Daswani, and M. Hutter. Self-modification of policy and utility function in rational agents. In \emph{Artificial General Intelligence (AGI 2016)}, LNCS 9782, pp. 1--11, 2016. \href{https://doi.org/10.1007/978-3-319-41649-6_1}{doi:10.1007/978-3-319-41649-6\_1}.
\bibitem{lean4} L. de Moura and S. Ullrich. The Lean 4 theorem prover and programming language. In \emph{Automated Deduction -- CADE 28}, LNCS 12699, pp. 625--635, 2021. \href{https://doi.org/10.1007/978-3-030-79876-5_37}{doi:10.1007/978-3-030-79876-5\_37}.
\bibitem{crttartifact} A. Strand. \emph{CRTT v1.0: full architecture and reconciled verification evidence}. Сопроводительный архив, ноябрь 2026 г. \texttt{CRTT-v1.0-FULL-ARCHITECTURE-RECONCILED.zip}. Идентифицируется по SHA-256 в разделе~\ref{sec:artifact}.
\bibitem{strand110observer2026} A. Strand. \emph{110 OBSERVER: Observability of a System's Own Past}. Препринт, 2026. \href{https://zenodo.org/records/24435829}{Запись Zenodo 24435829}.
\bibitem{strand110memory2026} A. Strand. \emph{Essential Inputs and Memory after Coalescence in Rule 110}. Препринт, 2026. \href{https://zenodo.org/records/24435942}{Запись Zenodo 24435942}.
\bibitem{strandcontrol2026} A. Strand. \emph{Representation, Incremental Memory Value, and Budgeted Control in Finite Channels}. Препринт, 2026. \href{https://doi.org/10.5281/zenodo.24458425}{doi:10.5281/zenodo.24458425}; \href{https://zenodo.org/records/24458425}{Запись Zenodo 24458425}.
\end{thebibliography}
% END ORIGINAL: source/ru/references.tex
\end{document}
% END ORIGINAL: source/ru/main.tex
